\scanpage{33}
\mspage{15}
\textcircled{2} \uline{Le groupe} $S\widehat{\widehat{\mathcal E}}$.

\unsure{Par|Pour} le calcul, il s'explicite comme l'ens des systèmes
\[
(1)\qquad\left(\underset{\hat{\mathbb Z}^{*}}{\underset{\rotatebox{90}{$\in$}}{p}},\underset{\unsure{\mathfrak S_3|\mathcal S_3}}{\underset{\rotatebox{90}{$\in$}}{\sigma}},\underset{\hat\pi}{\underset{\rotatebox{90}{$\in$}}{g_0}},\underset{\hat\pi}{\underset{\rotatebox{90}{$\in$}}{g_1}},\underset{\hat\pi}{\underset{\rotatebox{90}{$\in$}}{g_\infty}}\right)=\tilde u
\]
satisfaisant la condition
\[
(2)\qquad\operatorname{int}(g_\infty)(l_{\sigma(\infty)}^p)\cdot\operatorname{int}(g_1)(l_{\sigma(1)}^p)\cdot\operatorname{int}(g_{\unsure{0|\infty}})\unsure{(}l_{\sigma(\unsure{0|\infty})}^p)=1
\]
et tels que les trois facteurs de ce produit soient un syst. de générateurs de $\hat\pi$ [avec la relation génératrice (2)]. Le morphisme $u$ défini par ce \deleted{\unsure{morph.}} système est
\[
(3)\qquad\left\{\begin{aligned}
u(l_0)&\unsure{\mapsto|=}\operatorname{int}(g_0)l_{\sigma(0)}^p\\
u(l_1)&=\operatorname{int}(g_1)l_{\sigma(1)}^p\\
u(l_\infty)&=\operatorname{int}(g_\infty)l_{\sigma(\infty)}^p
\end{aligned}\right.
\]
\inserted{ce qui explicite l'hom can}
\[
(4)\qquad S\widehat{\widehat{\mathcal E}}\xrightarrow{p}\widehat{\widehat{\mathcal E}}=\operatorname{Aut}_{\mathrm{lac}}(\hat\pi)
\]
L'homomorphisme canonique
\[
(5)\qquad i:\hat\pi\times\hat{\mathbb Z}^{\{0,1,\infty\}}\longrightarrow S\widehat{\widehat{\mathcal E}}
\]
s'explicite par
\[
(6)\qquad i(g;\alpha_0,\alpha_1,\alpha_\infty)=(1,1,gl_0^{\alpha_0},gl_1^{\alpha_1},gl_\infty^{\alpha_\infty})
\]
donc
\[
(7)\quad\deleted{\text{donc}}\quad pi(g;\alpha_0,\alpha_1,\alpha_\infty)=\operatorname{int}g.
\]
\margin{On retrouve que \unsure{(4)|(p)} est surjectif, et a comme noyau $i(\hat\pi)$}
La multiplication dans $\widehat{\widehat{\mathcal E}}$ est donnée par
\[
(8)\quad\underbrace{(p,\sigma,g_0,g_1,g_\infty)}_{\tilde u}\underbrace{(p',\sigma',g'_0,g'_1,g'_\infty)}_{\tilde u'}=(pp',\sigma\sigma',u(g'_0)g_{\sigma'(0)},u(g'_1)g_{\sigma'(1)},u(g'_\infty)g_{\sigma'(\infty)}\unsure{)}
\]
Cette formule n'est pas, du point de vue du calcul
\mspage{32}
